\documentclass{beamer}

\title{Low-cost perfusion pump}
\subtitle{Multidisciplinary design (GM1)}
\author{L. Buxton \and L. McKiernan \and I. Henry \and L. Payne \and T. Azam}
\institute{University of Cambridge}
\date{March 2026}

\begin{document}

\frame{\titlepage}

\frame{
    \frametitle{Motivation}
}



\frame{
    \frametitle{Software}

    \begin{itemize}
        \item Customise the target pressure profile
              \begin{itemize}
                  \item via physiological parameters (e.g., heart rate, systolic/diastolic pressures)
                  \item via body state (e.g., age, sex, rest vs. exercise)
                  \item via vessel type (e.g., artery, capillary, vein)
              \end{itemize}
        \item View a live trace of the pressure sensor reading at the test start point
              \begin{itemize}
                  \item including the (cumulative) error (esp. during the syringe transition) and uptime
              \end{itemize}
        \item Set schedules for operation
              \begin{itemize}
                  \item e.g. ramping pressure profile to test vessel integrity
              \end{itemize}
        \item Flag extreme pressure values
              \begin{itemize}
                  \item e.g., loss of pressure from a sudden vessel rupture
                  \item e.g., increase in pressure from a blockage in the system from clotting
              \end{itemize}
        \item Export data traces as CSV
    \end{itemize}
}
\end{document}